Abbildungen können durch das Kommando \lstinline+\includegraphics+ aus dem
Paket~\ctanref{graphicx} eingebunden werden. Die Abbildung muss dafür als Grafikdatei
vorliegen. Legen Sie Abbildungen im Unterverzeichnis \texttt{figures} ab
(siehe Abschnitt~\ref{sec:tpl:dir}). Verwenden Sie je nach Anwendungsfall die
folgenden Bildformate.

\begin{description}
\item[PDF] Technische Zeichnungen, Grafiken, Diagramme, etc. -- Im Allgemeinen
  sind hier sog. Vektorgrafiken gemeint. Diese bestehen aus Punkten, Linien
  und Flächen und können ohne Qualitätsverlust vergrößert oder verkleinert
  werden. Zudem nehmen sie weniger Speicherplatz ein.

  Beachten Sie allerdings, dass Texte innerhalb von Vektorgrafiken dann
  ebenfalls mit skaliert werden und das u.\,U. zu einem uneinheitlichen
  Textbild führen kann.

\item[PNG] Screenshots, computergenerierte Abbildungen (3D-Visualisierungen,
  Simulationsergebnisse, \dots), Land- bzw. Wetterkarten, Heatmaps -- Hier
  handelt es sich um sog. Rastergrafiken. Solche Bilder bestehen aus einer
  schachbrettartigen Anordnung kleinster Pixel, die unterschiedlich gefärbt
  sind. Da die Pixel nicht einzeln wahrgenommen werden, entsteht ein
  Bildeindruck. Die Vergrößerung oder Verkleinerung von Rasterbildern führt
  allerdings zu Qualitätsverlusten (Unschärfe, Detailverlust). Rastergrafiken
  beanspruchen mehr Speicherplatz als Vektorgrafiken.

  Bei PNG-Dateien handelt es sich um ein verlustloses Bildformat, d.\,h. die
  Bildinformationen werden bei der Speicherung nicht verändert. Es eignet sich
  daher für Abbildungen mit feinen Details (z.\,B. Linien) und kontrastreiche
  Darstellungen (Schwarz-Weiß-Zeichnungen).

\item[JPEG] Fotografien, eingescannte Bildvorlagen -- Hierbei handelt es sich
  ebenfalls um Rastergrafiken. Es kommt jedoch ein verlustbehaftetes
  Speicherverfahren zum Einsatz. Farbwerte werden hierdurch verändert. Bei
  auf "`natürlichem Weg"' entstandenen Abbildungen wie Fotografien oder
  gescannten Dokumenten fallen diese Veränderungen nicht auf. Unnatürliche
  Abbildungen mit starken Kontrasten, feinen Linien oder stark gesättigten
  Farben führen jedoch zu Darstellungsfehlern (sog. Artefakten), weshalb das
  JPEG-Format in diesem Fall nicht verwendet werden sollte. Aufgrund der
  Speichertechnik benötigen JPEG-Bilder i.\,A. weniger Speicher als PNG-Bilder.
\end{description}

Abbildungen sollten immer innerhalb einer \texttt{figure}-Umgebung platziert
werden. Somit ist die Nummerierung der Abbildungen sichergestellt und eine
Aufnahme in das Abbildungsverzeichnis gesichert. Ein Beispiel findet sich in
Listing~\ref{lst:latex:figure}.

\begin{lstlisting}[language={[LaTeX]{TeX}},style=block,escapechar=\!,float,caption={Eine Abbildung},label={lst:latex:figure}]
\begin{figure}
\centering
\includegraphics[width=0.7\linewidth]{figures/!\lara{Dateiname}!}
\caption{!\lara{Bildunterschrift}!}
\label{!\lara{Verweismarke}!}
\end{figure}
\end{lstlisting}

Im Dateinamen muss -- und sollte -- keine Dateiendung angegeben werden. Die
Verweismarke kann entfallen, wenn kein Querverweis auf die Abbildung geplant
ist.

Bei der \texttt{figure}-Umgebung handelt es sich um eine Gleitumgebung. Die
Abbildung kann automatisch an eine andere Stelle des Dokuments verschoben
werden, wenn die Platzverhältnisse es erfordern. Weitere Hinweise hierzu finden
Sie in Abschnitt~\ref{sec:final:float}.
